\chapter{Proportions et évolutions}
\section*{Exercice 17.1 -- Une proportion}
$54$ personnes représentent $18\,\%$ d'un groupe. Quel est son effectif ?

\section*{Exercice 17.2 -- Proportion de proportion}
$40\,\%$ des élèves prennent un bus ; parmi eux, $30\,\%$ prennent la ligne A. Quelle proportion de tous les élèves prend cette ligne ?

\section*{Exercice 17.3 -- Successions}
Une quantité augmente de $10\,\%$, puis baisse de $15\,\%$. Donner le taux global.

\section*{Exercice 17.4 -- Réciproque}
Quel taux compense exactement une baisse de $30\,\%$ ?

\section*{Exercice 17.5 -- Points ou pourcentage}
Une proportion passe de $12\,\%$ à $15\,\%$. Donner la hausse en points et la hausse relative.

\section*{Exercice 17.6 -- Pondérer}
Un groupe de $40$ personnes compte $25\,\%$ de réponses oui ; un autre de $60$ en compte $50\,\%$. Calculer la proportion globale.
